\section{LTspice Guide}

This guide is written for students using LTspice for the first time. All screenshots are taken from
LTspice \ltspiceVersion.

\subsection{Step 1: Starting LTspice and creating a schematic}
\todo{File $\rightarrow$ New Schematic; save as \texttt{Lab1/Ex01.asc}.}

\guidefig{step01_start}{Starting LTspice and creating a schematic}

\subsection{Step 2: Placing components and wires}
\todo{Component (F2), Resistor (R), Wire (F3), Ground (G), Net label (F4); rotate with Ctrl+R.}

\guidefig{step02_place}{Placing components and wires}

\subsection{Step 3: Editing component values}
\todo{Right-click a component; note the suffixes \texttt{k}, \texttt{meg}, \texttt{m}, \texttt{u}.}

\guidefig{step03_values}{Editing component values}

\subsection{Step 4: Dependent sources}
\todo{Place \texttt{e} (VCVS) / \texttt{h} (CCVS) and set the gain.}

\guidefig{step04_dep_source}{Dependent sources}

\subsection{Step 5: Running a DC operating point analysis}
\todo{Simulate $\rightarrow$ Edit Simulation Cmd $\rightarrow$ DC op pnt; run.}

\guidefig{step05_op}{Running a DC operating point analysis}

\subsection{Step 6: Reading the results}
\todo{Hover over wires/components, place \texttt{.op} data labels, open the SPICE Error Log.}

\guidefig{step06_op_result}{Reading the results}

\subsection{Step 7: Automating measurements with \texttt{.meas}}
\todo{Use \texttt{.meas} to print values to the log for comparison.}

\guidefig{step07_meas}{Automating measurements with \texttt{.meas}}

\begin{tipbox}
\todo{Add one practical tip for this lab (common error, shortcut, or setting).}
\end{tipbox}
